% jcp-header.tex (jcpquarto 1.0.0)
% Preamble for Journal of Consumer Psychology (JCP) manuscripts. pdfLaTeX fonts: Times for the text,
% tables, captions, and notes; Helvetica for any sans serif text.

% ---- Fonts ----
\usepackage{mathptmx}
\usepackage[scaled=0.92]{helvet}

% ---- Quarto 1.4 (pandoc before 3.2) ----
% Those versions wrap \paragraph and \subparagraph in a way that titlesec
% cannot read ("Argument of \paragraph has an extra }"). The original
% commands are put back before titlesec is loaded; the heading formats below
% define them anew.
\makeatletter
\@ifundefined{xxxParagraphStar}{%
  \@ifundefined{oldparagraph}{}{\let\paragraph\oldparagraph}%
  \@ifundefined{oldsubparagraph}{}{\let\subparagraph\oldsubparagraph}%
}{}
\makeatother

% ---- Packages ----
\usepackage{float}
\usepackage{setspace}
\usepackage{ragged2e}
\usepackage{titlesec}
\usepackage{etoolbox}
\usepackage{longtable}
\usepackage{booktabs}
\usepackage[bottom]{footmisc}
\usepackage{indentfirst}
\usepackage{fancyhdr}

% ---- Characters typed directly ----
% pdfLaTeX has no glyphs of its own for Greek letters and for many signs that
% are typed as characters (a Greek alpha, a minus sign, "less than or equal").
% In the text, Greek letters are taken upright from the Symbol font, as they
% appear in a Word file; in formulas they are the usual math letters. The
% other signs are set in math mode. XeLaTeX and LuaLaTeX read all of these
% characters without help, so the definitions apply to pdfLaTeX only.
\usepackage{iftex}
\usepackage{pifont}
\ifPDFTeX
\providecommand{\texorpdfstring}[2]{#1}
\newcommand{\jcpgreek}[3]{\texorpdfstring{\TextOrMath{\Pisymbol{psy}{#1}}{#2}}{#3}}
\newcommand{\jcpsign}[1]{\texorpdfstring{\ensuremath{#1}}{}}
\DeclareUnicodeCharacter{03B1}{\jcpgreek{97}{\alpha}{alpha}}
\DeclareUnicodeCharacter{03B2}{\jcpgreek{98}{\beta}{beta}}
\DeclareUnicodeCharacter{03B3}{\jcpgreek{103}{\gamma}{gamma}}
\DeclareUnicodeCharacter{03B4}{\jcpgreek{100}{\delta}{delta}}
\DeclareUnicodeCharacter{03B5}{\jcpgreek{101}{\varepsilon}{epsilon}}
\DeclareUnicodeCharacter{03F5}{\jcpgreek{101}{\epsilon}{epsilon}}
\DeclareUnicodeCharacter{03B6}{\jcpgreek{122}{\zeta}{zeta}}
\DeclareUnicodeCharacter{03B7}{\jcpgreek{104}{\eta}{eta}}
\DeclareUnicodeCharacter{03B8}{\jcpgreek{113}{\theta}{theta}}
\DeclareUnicodeCharacter{03D1}{\jcpgreek{74}{\vartheta}{theta}}
\DeclareUnicodeCharacter{03B9}{\jcpgreek{105}{\iota}{iota}}
\DeclareUnicodeCharacter{03BA}{\jcpgreek{107}{\kappa}{kappa}}
\DeclareUnicodeCharacter{03BB}{\jcpgreek{108}{\lambda}{lambda}}
\DeclareUnicodeCharacter{03BC}{\jcpgreek{109}{\mu}{mu}}
\DeclareUnicodeCharacter{03BD}{\jcpgreek{110}{\nu}{nu}}
\DeclareUnicodeCharacter{03BE}{\jcpgreek{120}{\xi}{xi}}
\DeclareUnicodeCharacter{03BF}{\jcpgreek{111}{o}{o}}
\DeclareUnicodeCharacter{03C0}{\jcpgreek{112}{\pi}{pi}}
\DeclareUnicodeCharacter{03D6}{\jcpgreek{118}{\varpi}{pi}}
\DeclareUnicodeCharacter{03C1}{\jcpgreek{114}{\rho}{rho}}
\DeclareUnicodeCharacter{03C2}{\jcpgreek{86}{\varsigma}{sigma}}
\DeclareUnicodeCharacter{03C3}{\jcpgreek{115}{\sigma}{sigma}}
\DeclareUnicodeCharacter{03C4}{\jcpgreek{116}{\tau}{tau}}
\DeclareUnicodeCharacter{03C5}{\jcpgreek{117}{\upsilon}{upsilon}}
\DeclareUnicodeCharacter{03C6}{\jcpgreek{106}{\varphi}{phi}}
\DeclareUnicodeCharacter{03D5}{\jcpgreek{102}{\phi}{phi}}
\DeclareUnicodeCharacter{03C7}{\jcpgreek{99}{\chi}{chi}}
\DeclareUnicodeCharacter{03C8}{\jcpgreek{121}{\psi}{psi}}
\DeclareUnicodeCharacter{03C9}{\jcpgreek{119}{\omega}{omega}}
\DeclareUnicodeCharacter{0391}{\jcpgreek{65}{\mathrm{A}}{A}}
\DeclareUnicodeCharacter{0392}{\jcpgreek{66}{\mathrm{B}}{B}}
\DeclareUnicodeCharacter{0393}{\jcpgreek{71}{\Gamma}{Gamma}}
\DeclareUnicodeCharacter{0394}{\jcpgreek{68}{\Delta}{Delta}}
\DeclareUnicodeCharacter{2206}{\jcpgreek{68}{\Delta}{Delta}}
\DeclareUnicodeCharacter{0395}{\jcpgreek{69}{\mathrm{E}}{E}}
\DeclareUnicodeCharacter{0396}{\jcpgreek{90}{\mathrm{Z}}{Z}}
\DeclareUnicodeCharacter{0397}{\jcpgreek{72}{\mathrm{H}}{H}}
\DeclareUnicodeCharacter{0398}{\jcpgreek{81}{\Theta}{Theta}}
\DeclareUnicodeCharacter{0399}{\jcpgreek{73}{\mathrm{I}}{I}}
\DeclareUnicodeCharacter{039A}{\jcpgreek{75}{\mathrm{K}}{K}}
\DeclareUnicodeCharacter{039B}{\jcpgreek{76}{\Lambda}{Lambda}}
\DeclareUnicodeCharacter{039C}{\jcpgreek{77}{\mathrm{M}}{M}}
\DeclareUnicodeCharacter{039D}{\jcpgreek{78}{\mathrm{N}}{N}}
\DeclareUnicodeCharacter{039E}{\jcpgreek{88}{\Xi}{Xi}}
\DeclareUnicodeCharacter{039F}{\jcpgreek{79}{\mathrm{O}}{O}}
\DeclareUnicodeCharacter{03A0}{\jcpgreek{80}{\Pi}{Pi}}
\DeclareUnicodeCharacter{03A1}{\jcpgreek{82}{\mathrm{P}}{P}}
\DeclareUnicodeCharacter{03A3}{\jcpgreek{83}{\Sigma}{Sigma}}
\DeclareUnicodeCharacter{03A4}{\jcpgreek{84}{\mathrm{T}}{T}}
\DeclareUnicodeCharacter{03A5}{\jcpgreek{85}{\Upsilon}{Upsilon}}
\DeclareUnicodeCharacter{03A6}{\jcpgreek{70}{\Phi}{Phi}}
\DeclareUnicodeCharacter{03A7}{\jcpgreek{67}{\mathrm{X}}{X}}
\DeclareUnicodeCharacter{03A8}{\jcpgreek{89}{\Psi}{Psi}}
\DeclareUnicodeCharacter{03A9}{\jcpgreek{87}{\Omega}{Omega}}
\DeclareUnicodeCharacter{2212}{\texorpdfstring{\ensuremath{-}}{-}}
\DeclareUnicodeCharacter{2264}{\jcpsign{\leq}}
\DeclareUnicodeCharacter{2265}{\jcpsign{\geq}}
\DeclareUnicodeCharacter{2260}{\jcpsign{\neq}}
\DeclareUnicodeCharacter{2248}{\jcpsign{\approx}}
\DeclareUnicodeCharacter{2261}{\jcpsign{\equiv}}
\DeclareUnicodeCharacter{223C}{\jcpsign{\sim}}
\DeclareUnicodeCharacter{2245}{\jcpsign{\cong}}
\DeclareUnicodeCharacter{221D}{\jcpsign{\propto}}
\DeclareUnicodeCharacter{2213}{\jcpsign{\mp}}
\DeclareUnicodeCharacter{2217}{\jcpsign{\ast}}
\DeclareUnicodeCharacter{22C5}{\jcpsign{\cdot}}
\DeclareUnicodeCharacter{2219}{\jcpsign{\bullet}}
\DeclareUnicodeCharacter{221A}{\jcpsign{\surd}}
\DeclareUnicodeCharacter{221E}{\jcpsign{\infty}}
\DeclareUnicodeCharacter{2211}{\jcpsign{\sum}}
\DeclareUnicodeCharacter{220F}{\jcpsign{\prod}}
\DeclareUnicodeCharacter{222B}{\jcpsign{\int}}
\DeclareUnicodeCharacter{2202}{\jcpsign{\partial}}
\DeclareUnicodeCharacter{2207}{\jcpsign{\nabla}}
\DeclareUnicodeCharacter{2208}{\jcpsign{\in}}
\DeclareUnicodeCharacter{2209}{\jcpsign{\notin}}
\DeclareUnicodeCharacter{2282}{\jcpsign{\subset}}
\DeclareUnicodeCharacter{2286}{\jcpsign{\subseteq}}
\DeclareUnicodeCharacter{2229}{\jcpsign{\cap}}
\DeclareUnicodeCharacter{222A}{\jcpsign{\cup}}
\DeclareUnicodeCharacter{2205}{\jcpsign{\emptyset}}
\DeclareUnicodeCharacter{2200}{\jcpsign{\forall}}
\DeclareUnicodeCharacter{2203}{\jcpsign{\exists}}
\DeclareUnicodeCharacter{2227}{\jcpsign{\wedge}}
\DeclareUnicodeCharacter{2228}{\jcpsign{\vee}}
\DeclareUnicodeCharacter{2032}{\jcpsign{{}^{\prime}}}
\DeclareUnicodeCharacter{2033}{\jcpsign{{}^{\prime\prime}}}
\DeclareUnicodeCharacter{2194}{\jcpsign{\leftrightarrow}}
\DeclareUnicodeCharacter{21D2}{\jcpsign{\Rightarrow}}
\DeclareUnicodeCharacter{21D0}{\jcpsign{\Leftarrow}}
\DeclareUnicodeCharacter{21D4}{\jcpsign{\Leftrightarrow}}
\DeclareUnicodeCharacter{2197}{\jcpsign{\nearrow}}
\DeclareUnicodeCharacter{2198}{\jcpsign{\searrow}}
\DeclareUnicodeCharacter{25CF}{\jcpsign{\bullet}}
\DeclareUnicodeCharacter{25CB}{\jcpsign{\circ}}
\DeclareUnicodeCharacter{25A0}{\jcpsign{\blacksquare}}
\DeclareUnicodeCharacter{25A1}{\jcpsign{\square}}
\DeclareUnicodeCharacter{25B2}{\jcpsign{\blacktriangle}}
\DeclareUnicodeCharacter{25B3}{\jcpsign{\triangle}}
\DeclareUnicodeCharacter{25BC}{\jcpsign{\blacktriangledown}}
\DeclareUnicodeCharacter{2605}{\jcpsign{\bigstar}}
\DeclareUnicodeCharacter{2606}{\texorpdfstring{\ding{73}}{}}
\DeclareUnicodeCharacter{2023}{\jcpsign{\triangleright}}
\DeclareUnicodeCharacter{2713}{\texorpdfstring{\ding{51}}{}}
\DeclareUnicodeCharacter{2714}{\texorpdfstring{\ding{52}}{}}
\DeclareUnicodeCharacter{2717}{\texorpdfstring{\ding{55}}{}}
\DeclareUnicodeCharacter{2718}{\texorpdfstring{\ding{56}}{}}
\DeclareUnicodeCharacter{2070}{\textsuperscript{0}}
\DeclareUnicodeCharacter{2074}{\textsuperscript{4}}
\DeclareUnicodeCharacter{2075}{\textsuperscript{5}}
\DeclareUnicodeCharacter{2076}{\textsuperscript{6}}
\DeclareUnicodeCharacter{2077}{\textsuperscript{7}}
\DeclareUnicodeCharacter{2078}{\textsuperscript{8}}
\DeclareUnicodeCharacter{2079}{\textsuperscript{9}}
\DeclareUnicodeCharacter{207A}{\textsuperscript{+}}
\DeclareUnicodeCharacter{207B}{\textsuperscript{\ensuremath{-}}}
\DeclareUnicodeCharacter{2080}{\textsubscript{0}}
\DeclareUnicodeCharacter{2081}{\textsubscript{1}}
\DeclareUnicodeCharacter{2082}{\textsubscript{2}}
\DeclareUnicodeCharacter{2083}{\textsubscript{3}}
\DeclareUnicodeCharacter{2084}{\textsubscript{4}}
\DeclareUnicodeCharacter{2085}{\textsubscript{5}}
\DeclareUnicodeCharacter{2086}{\textsubscript{6}}
\DeclareUnicodeCharacter{2087}{\textsubscript{7}}
\DeclareUnicodeCharacter{2088}{\textsubscript{8}}
\DeclareUnicodeCharacter{2089}{\textsubscript{9}}
\DeclareUnicodeCharacter{2153}{\textsuperscript{1}/\textsubscript{3}}
\DeclareUnicodeCharacter{2154}{\textsuperscript{2}/\textsubscript{3}}
\DeclareUnicodeCharacter{2060}{\nobreak}
\DeclareUnicodeCharacter{FB00}{ff}
\DeclareUnicodeCharacter{FB01}{fi}
\DeclareUnicodeCharacter{FB02}{fl}
\DeclareUnicodeCharacter{FB03}{ffi}
\DeclareUnicodeCharacter{FB04}{ffl}
\fi

% ---- Tables: single spaced, in the text font ----
% \setstretch{1} rather than \singlespacing, which adds a blank line
\AtBeginEnvironment{table}{\normalsize\setstretch{1}}
\AtBeginEnvironment{table*}{\normalsize\setstretch{1}}
\AtBeginEnvironment{longtable}{\normalsize\setstretch{1}}
% Space above a table title, and less space between a table and its note
\setlength{\LTpre}{18pt}
\setlength{\LTpost}{6pt}
% A table that fits on a page is not split: \jcpkeep{height} starts a new
% page when less than that height is left on the current page
\newcommand{\jcpkeep}[1]{%
  \par\ifdim\pagetotal>0pt\ifdim\dimexpr\pagegoal-\pagetotal\relax<#1\relax\newpage\fi\fi}
% Tables and their titles are flush left
\setlength{\LTleft}{0pt}
\setlength{\LTright}{\fill}
% A table or a figure without a label takes no number. LaTeX counts every
% title, and pandoc 3.2 (in Quarto 1.5) counts every table, with or without
% one; the counter is put back afterward, so that the numbers agree with the
% Word file, and a title is printed without "Table 1" before it.
\newcounter{jcpsaved}
\newcounter{jcpplain}
\newcommand{\jcpplainbegin}[1]{\setcounter{jcpsaved}{\value{#1}}\stepcounter{jcpplain}%
  \begingroup\ifdefined\captionsetup\captionsetup{format=plain,labelformat=empty,labelsep=none}\fi%
  \expandafter\def\csname theH#1\endcsname{plain.\arabic{jcpplain}}}
\newcommand{\jcpplainend}[1]{\endgroup\setcounter{#1}{\value{jcpsaved}}}

% ---- Captions ----
% "Table 1" on the first line and the title on the second, above the table
% and flush left; "Fig. 1. Caption." on one line. APA sets the table number
% in bold (7th edition) and the title in italics.
\usepackage[font={rm,normalsize,stretch=2},justification=raggedright,
            singlelinecheck=false,skip=0pt]{caption}
\captionsetup[table]{position=above,labelsep=newline,labelfont={bf},textfont={it}}
\captionsetup[figure]{position=below,labelsep=period,labelfont={},textfont={}}
% The captions of the panels of a figure: "(a) Text", centered, in the type
% of the text. (Quarto loads the subcaption package when a document has a
% figure with panels.)
\makeatletter
\AtBeginDocument{\@ifpackageloaded{subcaption}{%
  \captionsetup[sub]{font={rm,normalsize,singlespacing},labelfont={},textfont={},
    labelformat=parens,labelsep=space,justification=centering,singlelinecheck=false,skip=4pt}}{}}
\makeatother
% The caption of a picture in a panel without a number: the same, without a letter
\newcommand{\jcppanelcaption}[1]{{\captionsetup{font={rm,normalsize,singlespacing},labelfont={},
    textfont={},justification=centering,singlelinecheck=false,skip=4pt}\caption*{#1}}}

% ---- No underlining ----
\renewcommand{\underline}[1]{#1}

% ---- Notes under tables and figures ----
% Written in the .qmd file as a ::: {.notes} div; the filter adds the label
\newenvironment{jcpnotes}{%
  \par\nopagebreak[4]\vskip 4pt%
  \begingroup\normalsize\setstretch{1}\raggedright
  \setlength{\parindent}{0pt}\setlength{\parskip}{2pt}%
}{%
  \par\endgroup\vskip 6pt%
}

% ---- Hypotheses: indented, with 12 pt of space above and below ----
% Written in the .qmd file as a ::: {.hypothesis} div; the filter sets the
% label ("H1:") in bold. \addvspace keeps the space single between two
% hypotheses in a row.
\newenvironment{jcphypothesis}{%
  \par\addvspace{12pt}\begingroup\setlength{\parindent}{0pt}%
  \setlength{\leftskip}{0.5in}%
}{%
  \par\endgroup\addvspace{12pt}%
}

% ---- Centered text, unindented text, and "insert table here" lines ----
\newenvironment{jcpcentered}{\par\begingroup\centering\setlength{\parindent}{0pt}}{\par\endgroup}
\newenvironment{jcpflush}{\par\begingroup\setlength{\parindent}{0pt}}{\par\endgroup}
\newenvironment{jcpcallout}{%
  \par\addvspace{0.5\baselineskip}\begingroup\centering\setlength{\parindent}{0pt}%
}{%
  \par\endgroup\addvspace{0.5\baselineskip}%
}

% ---- References: half-inch hanging indent ----
\ifdefined\cslhangindent
  \setlength{\cslhangindent}{0.5in}
\else
  \newlength{\cslhangindent}
  \setlength{\cslhangindent}{0.5in}
\fi
\newenvironment{jcpreferences}{\par\begingroup}{\par\endgroup}

% ---- Statistics helpers ----
\newcommand{\Mstat}[1]{\ensuremath{M = #1}}
\newcommand{\SDstat}[1]{\ensuremath{\text{SD} = #1}}
\newcommand{\pval}[1]{\ensuremath{p = #1}}
\newcommand{\Fstat}[3]{\ensuremath{F(#1, #2) = #3}}
\newcommand{\tstat}[2]{\ensuremath{t(#1) = #2}}
\newcommand{\chisq}[2]{\ensuremath{\chi^2(#1) = #2}}
\newcommand{\CI}[2]{95\% CI [#1, #2]}

% ---- Front matter (the filter places the text in these environments) ----
% Room for the page number in the header
\setlength{\headheight}{15pt}
% Bold labels (7th edition); the 6th edition sets them in plain type
\newcommand{\jcplabelfont}{\bfseries}
% Title, "Abstract", and "Author Note": centered
\newenvironment{jcptitle}{\par\begingroup\centering\jcplabelfont\setlength{\parindent}{0pt}}{\par\endgroup}
\newenvironment{jcpfronthead}{\par\begingroup\centering\jcplabelfont\setlength{\parindent}{0pt}}{\par\endgroup}
% Abstract: one paragraph without indentation. Keywords: indented.
\newenvironment{jcpabstract}{\par\begingroup\setlength{\parindent}{0pt}}{\par\endgroup}
\newenvironment{jcpkeywords}{\par\begingroup\setlength{\parindent}{0.5in}}{\par\endgroup}

% ---- MANUSCRIPT SETUP (submission format) ----
% APA manuscript: double spaced (linestretch), not justified or hyphenated,
% the page number at the top right. Headings (7th edition): level 1 centered
% and bold; level 2 flush left and bold; level 3 flush left, bold italic;
% levels 4 and 5 indented, run into the text, with a period.
\newcommand{\jcpsetupmanuscript}{%
  \pagestyle{fancy}\fancyhf{}\renewcommand{\headrulewidth}{0pt}\fancyhead[R]{\thepage}%
  \fancypagestyle{plain}{\fancyhf{}\renewcommand{\headrulewidth}{0pt}\fancyhead[R]{\thepage}}%
  \raggedright\hyphenpenalty=10000\exhyphenpenalty=10000%
  \setlength{\parindent}{0.5in}%
  \titleformat{\section}%
    {\normalfont\normalsize\bfseries\centering}{}{0pt}{}%
  \titlespacing*{\section}{0pt}{0pt}{0pt}%
  \titleformat{\subsection}%
    {\normalfont\normalsize\bfseries\raggedright}{}{0pt}{}%
  \titlespacing*{\subsection}{0pt}{0pt}{0pt}%
  \titleformat{\subsubsection}%
    {\normalfont\normalsize\bfseries\itshape\raggedright}{}{0pt}{}%
  \titlespacing*{\subsubsection}{0pt}{0pt}{0pt}%
  \titleformat{\paragraph}[runin]%
    {\normalfont\normalsize\bfseries}{}{0pt}{}[.]%
  \titlespacing*{\paragraph}{\parindent}{0pt}{0.5em}%
  \titleformat{\subparagraph}[runin]%
    {\normalfont\normalsize\bfseries\itshape}{}{0pt}{}[.]%
  \titlespacing*{\subparagraph}{\parindent}{0pt}{0.5em}%
}
% 6th edition: levels 3 to 5 run into the text; labels in plain type
\newcommand{\jcpsixthedition}{%
  \renewcommand{\jcplabelfont}{}%
  \titleformat{\subsubsection}[runin]%
    {\normalfont\normalsize\bfseries}{}{0pt}{}[.]%
  \titlespacing*{\subsubsection}{\parindent}{0pt}{0.5em}%
  \titleformat{\paragraph}[runin]%
    {\normalfont\normalsize\bfseries\itshape}{}{0pt}{}[.]%
  \titleformat{\subparagraph}[runin]%
    {\normalfont\normalsize\itshape}{}{0pt}{}[.]%
  \captionsetup[table]{labelfont={}}%
}

% ---- ARTICLE SETUP (publication-style preview; not for submission) ----
\newcommand{\jcpsetuparticle}{%
  \pagestyle{fancy}%
  \fancyhf{}%
  \fancyfoot[C]{\small\thepage}%
  \renewcommand{\headrulewidth}{0pt}%
  \fancypagestyle{plain}{%
    \fancyhf{}%
    \fancyfoot[C]{\small\thepage}%
    \renewcommand{\headrulewidth}{0pt}%
  }%
  \thispagestyle{plain}%
  \setlength{\parindent}{1.5em}%
  \titleformat{\section}%
    {\normalfont\large\bfseries\centering}{}{0pt}{}%
  \titlespacing*{\section}{0pt}{18pt}{6pt}%
  \titleformat{\subsection}%
    {\normalfont\normalsize\bfseries\raggedright}{}{0pt}{}%
  \titlespacing*{\subsection}{0pt}{14pt}{4pt}%
  \titleformat{\subsubsection}%
    {\normalfont\normalsize\bfseries\itshape\raggedright}{}{0pt}{}%
  \titlespacing*{\subsubsection}{0pt}{10pt}{3pt}%
  \titleformat{\paragraph}[runin]%
    {\normalfont\normalsize\bfseries}{}{0pt}{}[.]%
  \titlespacing*{\paragraph}{\parindent}{8pt}{0.5em}%
  \titleformat{\subparagraph}[runin]%
    {\normalfont\normalsize\bfseries\itshape}{}{0pt}{}[.]%
  \titlespacing*{\subparagraph}{\parindent}{8pt}{0.5em}%
  \captionsetup{font={rm,small,singlespacing},skip=6pt}%
  \captionsetup{font=small}%
}
\newenvironment{jcparttitle}{%
  \par\begingroup\fontsize{16}{19}\selectfont\bfseries\raggedright\setlength{\parindent}{0pt}%
}{\par\endgroup\vskip 10pt}
\newenvironment{jcpartauthors}{%
  \par\begingroup\centering\normalsize\bfseries\setlength{\parindent}{0pt}%
}{\par\endgroup\vskip 6pt}
\newenvironment{jcpartabstract}{%
  \par\noindent\textbf{Abstract}\par\vskip 4pt\begingroup\small\setlength{\parindent}{0pt}%
}{\par\endgroup\vskip 8pt}
\newenvironment{jcpartkeywords}{%
  \par\begingroup\small\setlength{\parindent}{0pt}%
}{\par\endgroup\vskip 8pt}
\newcommand{\jcpartrule}{%
  \vskip 0.8\baselineskip\noindent\rule{\textwidth}{0.4pt}\vskip 0.8\baselineskip}
\makeatletter
\newcommand{\jcpauthornote}[1]{%
  \begingroup\def\@makefnmark{}%
  \footnotetext{\fontsize{8}{10}\selectfont\raggedright #1}\endgroup}
\makeatother
